\documentclass[12pt,letterpaper]{article}
\usepackage{fullpage}
\usepackage[top=2cm, bottom=4.5cm, left=2.5cm, right=2.5cm]{geometry}
\usepackage{amsmath,amsthm,amsfonts,amssymb,amscd}
\usepackage{lastpage}
\usepackage{enumerate}
\usepackage{fancyhdr}
\usepackage{mathrsfs}
\usepackage{xcolor}
\usepackage{graphicx}
\usepackage{listings}
\usepackage{hyperref}

\hypersetup{%
  colorlinks=true,
  linkcolor=blue,
  linkbordercolor={0 0 1}
}
 
\renewcommand\lstlistingname{Algorithm}
\renewcommand\lstlistlistingname{Algorithms}
\def\lstlistingautorefname{Alg.}

\lstdefinestyle{Python}{
    language        = Python,
    frame           = lines, 
    basicstyle      = \footnotesize,
    keywordstyle    = \color{blue},
    stringstyle     = \color{green},
    commentstyle    = \color{red}\ttfamily
}

\setlength{\parindent}{0.0in}
\setlength{\parskip}{0.05in}

% Edit these as appropriate
\newcommand\course{CSE 3500}
\newcommand\hwnumber{1}                  % <-- homework number
\newcommand\NetIDa{netid19823}           % <-- NetID of person #1
\newcommand\NetIDb{netid12038}           % <-- NetID of person #2 (Comment this line out for problem sets)

\pagestyle{fancyplain}
\headheight 35pt
\lhead{MATH 20250, Prof. Lanier,
Midterm.}
                % <-- Comment this line out for problem sets (make sure you are person #1)


%%%%%%%%%%%%%%%%%%%%%%%%
\chead{\textbf{Name:}}
\rhead{}

\lfoot{}
\cfoot{}
\rfoot{\small\thepage}
\headsep .5em
\begin{document}
There are 100 points in the exam, over 9 numbered problems, some with multiple parts.\\
The total number of points for each problem is indicated.\\
Each part of a problem is worth the same amount.\\
There is also a bonus.


\vskip.1in
\hrule
\vskip.1in

1. (16 points) Give an example of each of the following, or explain why it is not possible.

\vskip.11in

a) Two non-square matrices that are column equivalent but not row equivalent.

\vskip1.5in

b) A real matrix $A$ with col($A$)=$\text{span}\Big(\Big\{\begin{pmatrix} 
-1\\0\\2\end{pmatrix},
\begin{pmatrix} 
4\\5\\6\end{pmatrix}\Big\}\Big)$ and nullspace($A$)= $\text{span}\Big(\Big\{\begin{pmatrix} 
1\\1\\-1\end{pmatrix}\Big\}\Big)$

\vskip1.5in
c) A real invertible matrix $A$ such that $\det(A^2)=-1$
\vskip1.5in

d) Two distinct subspaces $U,V$ of $\mathbb{R}^4$ so that their intersection is a subspace of dimension 1.

\newpage
\vskip.5in
2. (8 points) 

\vskip.1in

a) Give a precise definition for the following linear algebra term.
\vskip.15in
{\bf linear transformation:}


\vskip3in

b) Give an example of a matrix that is in row echelon form, but not in reduced row echelon form. Give the reasons why your matrix satisfies the definition of REF, but not of RREF.




\newpage

3. (10 points) Complete the following:

\vskip.15in

a) State the Rank-Nullity Theorem. Give an example that illustrates it.

\vskip3.5in

b) Give three examples of conditions that are equivalent to the following: the rows of an $n \times n$ matrix $A$ are linearly dependent.



\newpage

\vskip.3in

4. (12 points) Prove \textbf{one} of the following:

a) If $A$ and $A^{-1}$ are real invertible matrices with all integer entries, then either both matrices have determinant $1$, or both matrices have determinant $-1$.

b) If a homogeneous system of linear equations over $\mathbb{R}$ has more variables than equations, then the system has infinitely many solutions.

c) Let $V$ be a vector space over a field $\mathbb{F}$. Let $c$ be an element of $\mathbb{F}$ and let $\mathbf{v}$ be an element
of $V$. Suppose that $c\mathbf{v} = 0$. Use the defining properties of a vector space to show that either $c = 0$ or $\mathbf{v} = \mathbf{0}$.

\newpage
\vskip.15in

5. (25 points) Complete \textbf{all} of the following. Always, sometimes, or never? If it is always or never true, give a proof. If it is sometimes true, give an example where it’s true and where it’s false.

\vskip.15in

a) For square matrices $A$ and B, $\det(A+B)=\det(A)+\det(B)$.

\vskip2.5in

b) Two matrices that commute are row equivalent.

\vskip2.5in

c) If two distinct subspaces of $V$ have bases $B_1$ and $B_2$, then the set $B_1 \cap B_2$ is linearly independent.

\newpage

d) A linear transformation from $\mathbb{R}^6$ to $\mathbb{R}^8$ is one-to-one.

\vskip3.5in

e) Given a set of three distinct vectors in a vector space $V$, their span is a subspace of $V$.


\newpage



\vskip.3in
6. (5 points) Show how to express the determinant of the following matrix by applying cofactor expansion once, along any row or column of your choosing. Just set up the calculation, you do not need to further evaluate it.

$\begin{pmatrix}
    a&b&c&d\\
    e&f&g&h\\
    i&j&k&l\\
    m&n&o&p
\end{pmatrix}$

\vskip3in



7. (6 points) Consider the function $S: \mathcal{P}_3 \to \mathcal{P}_6$ that takes each real cubic polynomial of degree at most 3 to its square (that is, to the polynomial times itself). Is this a linear transformation? Explain why or why not.




\newpage



8. (12 points) Complete \textbf{one} of the following:

a) Compute a basis for each of the four fundamental subspaces of the real matrix $\begin{pmatrix}
    1&3&2\\
    0&0&1
\end{pmatrix}$.

b) Show that vector space $(\mathbb{Z}/3\mathbb{Z})^3$ has exactly 28 subspaces.







\newpage

9.  (6 points) Complete \textbf{one} of the following.

a) Suppose that for a linear transformation $T: \mathbb{R}^4 \to \mathbb{R}^3$ we have $T(\mathbf{u})=3\mathbf{w}$ and $T(\mathbf{v})=8\mathbf{w}$, where $\mathbf{w}$ is a nonzero vector in the codomain. Give two distinct examples of vectors that are in the kernel of $T$; or explain why this is not possible.

b) Give an example of a matrix that can be written as a product of exactly four elementary matrices in two different ways (and give the ways); or explain why this is not possible.
\newpage

\newpage

\vskip.5in

\textbf{Bonus:} Explain one or more linear algebra ideas or results that you've learned in the course so far but that haven't been asked about on this exam.
\end{document}
